\documentclass{beamer}

\usepackage[T2A]{fontenc}
\usepackage[utf8]{inputenc}
\usepackage[english,russian]{babel}
\pdfmapfile{+sansmathaccent.map}


\mode<presentation>
{
	\usetheme{Warsaw} % or try Darmstadt, Madrid, Warsaw, Rochester, CambridgeUS, ...
	\usecolortheme{seahorse} % or try seahorse, beaver, crane, wolverine, ...
	\usefonttheme{serif}  % or try serif, structurebold, ...
	\setbeamertemplate{navigation symbols}{}
	\setbeamertemplate{caption}[numbered]
} 


%%%%%%%%%%%%%%%%%%%%%%%%%%%%
% itemize settings


%%%%%%%%%%%%%%%%%%%%%%%%%%%%
% itemize settings

\definecolor{myhotpink}{RGB}{255, 80, 200}
\definecolor{mywarmpink}{RGB}{255, 60, 160}
\definecolor{mylightpink}{RGB}{255, 80, 200}
\definecolor{mypink}{RGB}{255, 30, 80}
\definecolor{mydarkpink}{RGB}{155, 25, 60}

\definecolor{mypaleblue}{RGB}{240, 240, 255}
\definecolor{mylightblue}{RGB}{120, 150, 255}
\definecolor{myblue}{RGB}{90, 90, 255}
\definecolor{mygblue}{RGB}{70, 110, 240}
\definecolor{mycourseblue}{RGB}{120, 140, 255}
\definecolor{mycourselightblue}{RGB}{245, 250, 255}
\definecolor{mydarkblue}{RGB}{0, 0, 180}
\definecolor{myblackblue}{RGB}{40, 40, 120}

\definecolor{mygreen}{RGB}{0, 200, 0}
\definecolor{mygreen2}{RGB}{245, 255, 230}
\definecolor{mydarkgreen}{RGB}{0, 120, 0}
\definecolor{mydarkgreen}{RGB}{0, 120, 0}


\definecolor{mygray}{gray}{0.8}
\definecolor{mydarkgray}{RGB}{80, 80, 160}
\definecolor{mylightgray}{RGB}{160, 160, 160}

\definecolor{mydarkred}{RGB}{160, 30, 30}
\definecolor{mylightred}{RGB}{255, 150, 150}
\definecolor{myred}{RGB}{200, 110, 110}
\definecolor{myblackred}{RGB}{120, 40, 40}

\definecolor{mypink}{RGB}{255, 30, 80}
\definecolor{myhotpink}{RGB}{255, 80, 200}
\definecolor{mywarmpink}{RGB}{255, 60, 160}
\definecolor{mylightpink}{RGB}{255, 80, 200}
\definecolor{mydarkpink}{RGB}{155, 25, 60}
\definecolor{mywhitepink}{RGB}{255, 240, 240}

\definecolor{mydarkcolor}{RGB}{60, 25, 155}
\definecolor{mylightcolor}{RGB}{130, 180, 250}



\setbeamertemplate{itemize items}[default]

\setbeamertemplate{itemize item}{\color{myblackblue}$\blacksquare$}
\setbeamertemplate{itemize subitem}{\color{mygblue}$\blacktriangleright$}
\setbeamertemplate{itemize subsubitem}{\color{mygray}$\blacksquare$}

\setbeamercolor{palette quaternary}{fg=white,bg=mycourseblue}
\setbeamercolor{titlelike}{parent=palette quaternary}

\setbeamercolor{palette quaternary2}{fg=black,bg=mycourselightblue}
\setbeamercolor{frametitle}{parent=palette quaternary2}

\setbeamerfont{frametitle}{size=\Large,series=\scshape}
\setbeamerfont{framesubtitle}{size=\normalsize,series=\upshape}





%%%%%%%%%%%%%%%%%%%%%%%%%%%%
% block settings

\setbeamercolor{block title}{bg=red!30,fg=black}

\setbeamercolor*{block title example}{bg=mygreen!40!white,fg=black}

\setbeamercolor*{block body example}{fg= black, bg= mygreen2}


%%%%%%%%%%%%%%%%%%%%%%%%%%%%
% URL settings
\hypersetup{
	colorlinks=true,
	linkcolor=blue,
	filecolor=blue,      
	urlcolor=blue,
}

%%%%%%%%%%%%%%%%%%%%%%%%%%

\renewcommand{\familydefault}{\rmdefault}


\newcommand{\mycolor}[1] {
\setbeamercolor{normal text}{fg={#1}} % Recolors standard body text
\setbeamercolor{structure}{fg={#1}}   % Recolors bullets and structure items
\usebeamercolor[fg]{normal text}      % Applies the normal text color immediately
}


\usepackage{amsmath}
\usepackage{mathtools}

\usepackage{cancel}
\usepackage{centernot}

\usepackage{changepage}
\usepackage{subcaption}

\usepackage{qrcode}

\newcommand{\bo}[1] {\mathbf{#1}}
\newcommand{\R}{\mathbb{R}} 
\newcommand{\T}{^\top}     



\newcommand{\mydate}{Fall 2026}

\newcommand{\mygit}{\textcolor{blue}{\href{https://github.com/SergeiSa/Advanced-Control-2026}{github.com/SergeiSa/Advanced-Control-2026}}}

\newcommand{\myqr}{ \textcolor{black}{\qrcode[height=1.5in]{https://github.com/SergeiSa/Advanced-Control-2026}}
}




\newcommand{\bref}[2] {\textcolor{blue}{\href{#1}{#2}}}

%%%%%%%%%%%%%%%%%%%%%%%%%%%%
% code settings

\usepackage{listings}
\usepackage{color}
% \definecolor{mygreen}{rgb}{0,0.6,0}
% \definecolor{mygray}{rgb}{0.5,0.5,0.5}
\definecolor{mymauve}{rgb}{0.58,0,0.82}
\lstset{ 
	backgroundcolor=\color{white},   % choose the background color; you must add \usepackage{color} or \usepackage{xcolor}; should come as last argument
	basicstyle=\footnotesize,        % the size of the fonts that are used for the code
	breakatwhitespace=false,         % sets if automatic breaks should only happen at whitespace
	breaklines=true,                 % sets automatic line breaking
	captionpos=b,                    % sets the caption-position to bottom
	commentstyle=\color{mygreen},    % comment style
	deletekeywords={...},            % if you want to delete keywords from the given language
	escapeinside={\%*}{*)},          % if you want to add LaTeX within your code
	extendedchars=true,              % lets you use non-ASCII characters; for 8-bits encodings only, does not work with UTF-8
	firstnumber=0000,                % start line enumeration with line 0000
	frame=single,	                   % adds a frame around the code
	keepspaces=true,                 % keeps spaces in text, useful for keeping indentation of code (possibly needs columns=flexible)
	keywordstyle=\color{blue},       % keyword style
	language=Octave,                 % the language of the code
	morekeywords={*,...},            % if you want to add more keywords to the set
	numbers=left,                    % where to put the line-numbers; possible values are (none, left, right)
	numbersep=5pt,                   % how far the line-numbers are from the code
	numberstyle=\tiny\color{mygray}, % the style that is used for the line-numbers
	rulecolor=\color{black},         % if not set, the frame-color may be changed on line-breaks within not-black text (e.g. comments (green here))
	showspaces=false,                % show spaces everywhere adding particular underscores; it overrides 'showstringspaces'
	showstringspaces=false,          % underline spaces within strings only
	showtabs=false,                  % show tabs within strings adding particular underscores
	stepnumber=2,                    % the step between two line-numbers. If it's 1, each line will be numbered
	stringstyle=\color{mymauve},     % string literal style
	tabsize=2,	                   % sets default tabsize to 2 spaces
	title=\lstname                   % show the filename of files included with \lstinputlisting; also try caption instead of title
}


%%%%%%%%%%%%%%%%%%%%%%%%%%%%
% URL settings
\hypersetup{
	colorlinks=false,
	linkcolor=blue,
	filecolor=blue,      
	urlcolor=blue,
}

%%%%%%%%%%%%%%%%%%%%%%%%%%

%%%%%%%%%%%%%%%%%%%%%%%%%%%%
% tikz settings

\usepackage{tikz}
\tikzset{every picture/.style={line width=0.75pt}}
\usetikzlibrary{shapes.geometric, decorations.markings, positioning, patterns, calc, arrows.meta}

\newcommand{\myqrframe}{
	\begin{frame}
		\begin{center}
			\centerline{Слайды доступны на Github:}
			\bigskip
			\centerline{\mygit}
			\bigskip
			\myqr
		\end{center}
	\end{frame}
}


\subtitle{Современные методы управления роботами}
\author{Сергей Савин}
\date{2026}


\title{Методы градиента политики}
\begin{document}
	\maketitle
	
	\begin{frame}{Логарифмический трюк}
		\begin{flushleft}
			
			В этой лекции мы будем использовать два свойства логарифма.
			
			Первое:
			\[
			\frac{d}{dx}\log f(x)
			=
			\frac{1}{f(x)}\frac{df(x)}{dx}.
			\]
			
			Следовательно:
			\[
			\boxed{
				\frac{df(x)}{dx}
				=
				f(x)\frac{d}{dx}\log f(x)
			}
			\]
			
			То же тождество применимо к градиентам по
			вектору параметров.
			
		\end{flushleft}
	\end{frame}
	
	\begin{frame}{Траектории}
		\begin{flushleft}
			
			Дискретная траектория $\tau$ — это последовательность состояний
			и действий:
			
			\[
			\tau=(s_0,a_0,s_1,a_1,\ldots,s_{T-1},a_{T-1},s_T).
			\]
			
			В обучении мы получаем траекторию, запуская эпизод:
			
			\begin{itemize}
				\item Сэмплируем начальное состояние $s_0$.
				\item Сэмплируем действие $a_t$ согласно
				текущей политике $\pi_\theta(a_t|s_t)$.
				\item Позволяем среде сэмплировать следующее состояние
				$s_{t+1}$ согласно
				$P(s_{t+1}|s_t,a_t)$.
			\end{itemize}
			
			\vspace{0.2cm}
			
			\begin{itemize}
				\item $\pi_\theta(a_t|s_t)$ определяет вероятность
				действия при текущем состоянии.
				\item $P(s_{t+1}|s_t,a_t)$ определяет вероятность
				следующего состояния при текущем состоянии и действии.
			\end{itemize}
			
		\end{flushleft}
	\end{frame}
	
	\begin{frame}{Награда за траекторию}
		\begin{flushleft}
			
			Зная траекторию $\tau$, мы можем вычислить
			суммарную награду, связанную с ней:
			
			\[
			\boxed{
				R(\tau)=\sum_{t=0}^{T-1}r(s_t,a_t)
			}
			\]
			
			\begin{itemize}
				\item Это детерминированная величина для фиксированной
				траектории.
				\item Её легко вычислить по собранным данным:
				мы просто суммируем награды, наблюдаемые в течение эпизода.
			\end{itemize}
			
		\end{flushleft}
	\end{frame}
	
	\begin{frame}{Вероятность траектории}
		\begin{flushleft}
			
			Некоторые траектории более вероятны, чем другие,
			при заданной политике.
			
			Мы можем вычислить вероятность траектории как:
			
			\[
			\boxed{
				p_\theta(\tau)
				=
				p(s_0)
				\prod_{t=0}^{T-1}
				\pi_\theta(a_t|s_t)
				P(s_{t+1}|s_t,a_t)
			}
			\]
			
			где $p(s_0)$ — вероятность начального состояния.
			
			Каждый член описывает один шаг траектории:
			политика выбирает действие, а среда
			сэмплирует следующее состояние.
			
		\end{flushleft}
	\end{frame}
	
	\begin{frame}{Упрощение с помощью логарифма}
		\begin{flushleft}
			
			Вспомним, что логарифм превращает произведения в суммы:
			
			\[
			\log(fg)=\log f+\log g.
			\]
			
			Поэтому мы можем записать логарифм вероятности траектории как:
			
			\[
			\boxed{
				\begin{aligned}
					\log p_\theta(\tau)
					={}&\log p(s_0)\\
					&+\sum_{t=0}^{T-1}\log\pi_\theta(a_t|s_t)\\
					&+\sum_{t=0}^{T-1}
					\log P(s_{t+1}|s_t,a_t).
				\end{aligned}
			}
			\]
			
			Это выражение сделает дифференцирование намного проще.
			
		\end{flushleft}
	\end{frame}
	
	\begin{frame}{Дифференцирование вероятности траектории}
		\begin{flushleft}
			
			Используя логарифмический трюк:
			
			\[
			\nabla_\theta p_\theta(\tau)
			=
			p_\theta(\tau)
			\nabla_\theta\log p_\theta(\tau).
			\]
			
			Только члены политики явно зависят от $\theta$.
			Предполагается, что распределение начального состояния и динамика среды
			не зависят от параметров политики.
			
			Следовательно:
			
			\[
			\boxed{
				\nabla_\theta p_\theta(\tau)
				=
				p_\theta(\tau)
				\sum_{t=0}^{T-1}
				\nabla_\theta\log\pi_\theta(a_t|s_t)
			}
			\]
			
		\end{flushleft}
	\end{frame}
	
	\begin{frame}{REINFORCE (1/3)}
		\begin{flushleft}
			
			Мы определяем целевую функцию как ожидаемую награду:
			
			\[
			\boxed{
				J(\theta)
				=
				\sum_\tau p_\theta(\tau)R(\tau)
			}
			\]
			
			где суммирование ведётся по всем возможным траекториям.
			
			Поскольку награда фиксированной траектории не
			зависит явно от параметров политики:
			
			\[
			\nabla_\theta J(\theta)
			=
			\sum_\tau
			\nabla_\theta p_\theta(\tau)R(\tau).
			\]
			
		\end{flushleft}
	\end{frame}
	
	\begin{frame}{REINFORCE (2/3)}
		\begin{flushleft}
			
			Подставляя наше предыдущее тождество:
			
			\[
			\begin{aligned}
				\nabla_\theta J(\theta)
				=
				\sum_\tau
				&R(\tau) p_\theta(\tau)
				\sum_{t=0}^{T-1}
				\nabla_\theta\log\pi_\theta(a_t|s_t).
			\end{aligned}
			\]
			
			Эквивалентно, используя определение математического ожидания:
			
			\[
			\boxed{
				\nabla_\theta J(\theta)
				=
				\mathbb{E}_{\tau\sim p_\theta}
				\left[
				R(\tau)
				\sum_{t=0}^{T-1}
				\nabla_\theta\log\pi_\theta(a_t|s_t)
				\right]
			}
			\]
			
		\end{flushleft}
	\end{frame}
	
	\begin{frame}{REINFORCE (3/3)}
		\begin{flushleft}
			
			Мы можем оценить градиент, используя сэмплированную траекторию:
			
			\[
			\widehat{\nabla_\theta J}
			=
			R(\tau)
			\sum_{t=0}^{T-1}
			\nabla_\theta\log\pi_\theta(a_t|s_t).
			\]
			
			Затем мы можем обновить параметры политики, используя
			градиентный подъём:
			
			\[
			\boxed{
				\theta
				\leftarrow
				\theta+
				\alpha R(\tau)
				\sum_{t=0}^{T-1}
				\nabla_\theta\log\pi_\theta(a_t|s_t)
			}
			\]
			
			Это алгоритм REINFORCE.
			
		\end{flushleft}
	\end{frame}
	
	\begin{frame}{От награды за траекторию к награде-до-конца}
		\begin{flushleft}
			{\small 
			В настоящее время мы приписываем одну и ту же суммарную награду
			за траекторию $R(\tau)$ каждому действию в траектории.
			
			Рассмотрим действие, предпринятое в момент времени $t$.
			Награды, полученные до этого действия, не могли быть
			на него повлияны.
			
			Следовательно, мы можем заменить суммарную награду на
			награду-до-конца:
			
			\[
			\boxed{
				G_t=\sum_{k=t}^{T-1}r(s_k,a_k)
			}
			\]
			
			Выражение для градиента политики принимает вид:
			
			\[
			\boxed{
				\nabla_\theta J(\theta)
				=
				\mathbb{E}
				\left[
				\sum_{t=0}^{T-1}
				G_t\nabla_\theta\log\pi_\theta(a_t|s_t)
				\right]
			}
			\]
		}
		\end{flushleft}
	\end{frame}
	
	\begin{frame}{От награды-до-конца к преимуществу}
		\begin{flushleft}
			
			Функции ценности состояния и ценности действия:
			
			\[
			V^\pi(s_t)=\mathbb{E}[G_t|s_t],
			\]
			
			\[
			Q^\pi(s_t,a_t)=\mathbb{E}[G_t|s_t,a_t].
			\]
			
			Мы можем использовать их для определения преимущества:
			%
			\[
			\boxed{
				A^\pi(s_t,a_t)
				=
				Q^\pi(s_t,a_t)-V^\pi(s_t)
			}
			\]
			
			Преимущество измеряет, насколько лучше действие
			по сравнению со средним действием в данном состоянии.
			
			Мы можем заменить награду-до-конца на оценку преимущества
			без изменения ожидаемого градиента политики:
			%
			\[
			\boxed{
				\nabla_\theta J(\theta)
				=
				\mathbb{E}
				\left[
				\sum_{t=0}^{T-1}
				A^\pi(s_t,a_t)
				\nabla_\theta\log\pi_\theta(a_t|s_t)
				\right]
			}
			\]
			
		\end{flushleft}
	\end{frame}
	
	\begin{frame}{Отношение вероятностей}
		\begin{flushleft}
			
			Предположим, мы собрали данные, используя старую политику
			$\pi_{\theta_{\mathrm{old}}}$, и теперь оптимизируем
			новую политику $\pi_\theta$.
			
			Мы определяем отношение вероятностей:
			
			\[
			\boxed{
				r_t(\theta)
				=
				\frac{
					\pi_\theta(a_t|s_t)
				}{
					\pi_{\theta_{\mathrm{old}}}(a_t|s_t)
				}
			}
			\]
			
			Это отношение говорит нам, насколько новая политика изменяет
			вероятность действия по сравнению со старой политикой.
			
			\begin{itemize}
				\item $r_t(\theta)>1$: новая политика делает действие
				более вероятным.
				\item $r_t(\theta)<1$: новая политика делает действие
				менее вероятным.
				\item $r_t(\theta)=1$: вероятность не изменилась.
			\end{itemize}
			
		\end{flushleft}
	\end{frame}
	
	\begin{frame}{Отсечение в PPO}
		\begin{flushleft}
			
			Мы хотим предотвратить чрезмерно большие изменения политики.
			
			PPO использует отсечённую целевую функцию:
			
			\[
			\boxed{
				\begin{aligned}
					J^{\mathrm{CLIP}}(\theta)
					=
					\mathbb{E}_{\mathrm{old}}
					\Big[
					\min\big(&r_t(\theta)A_t,\\
					&\operatorname{clip}
					(r_t(\theta),1-\epsilon,1+\epsilon)A_t
					\big)
					\Big].
				\end{aligned}
			}
			\]
			
			\begin{itemize}
				\item Отсечение ограничивает стимул слишком сильно изменять
				вероятности действий.
				\item $\epsilon$ определяет интервал отсечения.
				\item Целевая функция максимизируется по $\theta$.
			\end{itemize}
			
		\end{flushleft}
	\end{frame}
\end{document}
	
	\myqrframe
	
	
	
	
	
	
\end{document}

